%% 5-9 グラフの面積は進んだ距離（第11回：積分の最初の章として書き足した。ページ数は自由・偶数）
\bkchapter{5}{9}{グラフの面積は進んだ距離}{\bklevel{2}}{グラフの面積は進んだ距離}
\bkmichishirube{2-3,4-5,5-5,5-8}{5-10,5-11,5-13}

\begin{goal}
\begin{itemize}[label={\color{bkTheme}●}, leftmargin=1.6\zw, itemsep=0pt, topsep=1pt]
\item \bklevel{2}\enspace 速さが変わるときも、「細かく分けて、かけて、たす」で進んだ距離を求められる。
\item \bklevel{2}\enspace 速度のグラフ（$v$-$t$ グラフ）を長方形・三角形・台形に分けて、進んだ距離を計算できる。
\item \bklevel{2}\enspace 速度がマイナスのとき、面積の符号の意味（変位と道のり）を説明できる。
\end{itemize}
\end{goal}

\begin{hatena}
\noindent\begin{minipage}[c]{0.58\linewidth}\raggedright
トラックには、走るあいだの速さを記録する機械（タコグラフ）がついています。記録されるのは速さだけです。

右は、ある日の記録を簡単にしたグラフです（横は時間 h、たては速さ km/h）。この記録だけから、\mbox{何km}走ったかがわかります。どうやって求めるのでしょう。
\end{minipage}\hfill
\begin{minipage}[c]{0.4\linewidth}\centering
\begin{tikzpicture}
\begin{axis}[bkplot, width=50mm, height=38mm, xmin=0, xmax=6.8, ymin=0, ymax=74, xtick={1,2,3,4,5,6}, ytick={20,40,60}, xlabel={$t$}, ylabel={$v$}]
  \addplot[bkcurve, line width=1.3pt] coordinates {(0,0) (0.5,40) (2,40) (2.5,0) (3,0) (3.5,60) (5.5,60) (6,0)};
\end{axis}
\end{tikzpicture}
\end{minipage}
\end{hatena}

2-3では、速さが一定のとき、「速さ×時間＝道のり」が、速さと時間のグラフの長方形の面積になることを学びました。でも、ほんとうの乗り物の速さは、いつも変わっています。

この章では、速さが変わるときにも、グラフの面積が進んだ距離になることを、1段ずつ確かめます。カギは「\textbf{細かく分けて、かけて、たす}」の3つの段です。これが、\textbf{積分}（せきぶん）の考え方の始まりです。

%% ---------------- やってみよう・思い出そう ----------------
\begin{yattemiyou}[方眼のマスを数えて距離を求めよう]
自転車の速さを、2秒ごとに記録しました。記録と記録のあいだは、速さがまっすぐ（一定の割合で）変わったとします。
\begin{bkcenter}
\begin{bktable}{colspec={Q[l,m,wd=26mm] *{6}{Q[c,m,wd=10mm]}}, rowsep=1.5pt}
時刻 $t$（s） & 0 & 2 & 4 & 6 & 8 & 10 \\
速さ $v$（m/s） & 0 & 2 & 4 & 4 & 4 & 2 \\
\end{bktable}
\end{bkcenter}
\begin{enumerate}[label={\bkcircled[bkGreen]{\arabic*}}, leftmargin=2.4\zw, itemsep=1pt, topsep=3pt]
\item 方眼紙に、横を時刻、たてを速さにして点をとり、となりの点を直線で結ぼう（$v$-$t$ グラフ）。1マスは、たて\qty{1}{m/s}、横\qty{1}{s}にしよう。
\item グラフと横軸のあいだにある方眼のマスを数えよう。ななめの線で半分に切られたマスは、2つで1マスと数えよう。
\item 1マスは、何mの距離にあたるかな？　10秒で何m進んだかな？
\end{enumerate}
\end{yattemiyou}

\noindent やってみようのグラフをかくと、下のようになります。

\begin{bkfigure}
\begin{tikzpicture}
\begin{axis}[bkplot, width=112mm, height=52mm, xmin=0, xmax=11.2, ymin=0, ymax=5.6, xtick={1,...,10}, ytick={1,...,5}, xlabel={$t$（s）}, ylabel={$v$（m/s）}]
  \fill[bkYellow!40] (axis cs:0,0) -- (axis cs:4,4) -- (axis cs:4,0) -- cycle;
  \fill[bkBlue!20] (axis cs:4,0) rectangle (axis cs:8,4);
  \fill[bkGreen!25] (axis cs:8,0) -- (axis cs:8,4) -- (axis cs:10,2) -- (axis cs:10,0) -- cycle;
  \addplot[bkcurve, line width=1.4pt, mark=*, mark size=1.4pt] coordinates {(0,0) (2,2) (4,4) (6,4) (8,4) (10,2)};
  \node[font=\sffamily\footnotesize, fill=bkYellow!40, inner sep=1pt] at (axis cs:3.05,0.5) {8マス};
  \node[font=\sffamily\footnotesize, fill=bkBlue!20, inner sep=1pt] at (axis cs:6,2) {16マス};
  \node[font=\sffamily\footnotesize, fill=bkGreen!25, inner sep=1pt] at (axis cs:9,1) {6マス};
\end{axis}
\end{tikzpicture}
\bkfigcaption{0〜4秒は三角形、4〜8秒は長方形、8〜10秒は台形。}
\end{bkfigure}

\noindent\begin{minipage}[c]{0.66\linewidth}
\textbf{1マスは何m？}\quad 1マスは、たてが\qty{1}{m/s}、横が\qty{1}{s}です。「1秒に\qty{1}{m}進む速さで、1秒走った」ということなので、進んだ距離は\qty{1}{m}です。だから、\textbf{1マスは\qty{1}{m}}にあたります。
\end{minipage}\hfill
\begin{minipage}[c]{0.3\linewidth}\centering
\begin{tikzpicture}[x=1mm, y=1mm]
  \fill[bkBlue!20] (0,0) rectangle (14,14);
  \draw[bkBlue!80!black, line width=0.8pt] (0,0) rectangle (14,14);
  \node[font=\sffamily\small] at (7,7) {\qty{1}{m}};
  \node[font=\sffamily\footnotesize, anchor=north] at (7,-1.5) {横 \qty{1}{s}};
  \node[font=\sffamily\footnotesize, anchor=east] at (-1.5,7) {たて \qty{1}{m/s}};
\end{tikzpicture}
\end{minipage}

\noindent\begin{minipage}[c]{0.62\linewidth}\raggedright
\textbf{マスを数える}\quad 0〜4秒の三角形を大きくしたのが右の図です。まるごとのマス（黄色）が $0+1+2+3=6$ こ、ななめの線で半分に切られたマス（オレンジ）が4こあります。半分のマスは2つで1マスなので、4こで2マスです。合わせて $6+2=8$ マスです。
\end{minipage}\hfill
\begin{minipage}[c]{0.34\linewidth}\centering
\begin{tikzpicture}[x=7mm, y=7mm]
  \foreach \k in {1,2,3} {\foreach \j in {1,...,\k} {\fill[bkYellow!55] (\k,\j-1) rectangle (\k+1,\j);}}
  \foreach \k in {0,1,2,3} {\fill[bkOrange!40] (\k,\k) -- (\k+1,\k) -- (\k+1,\k+1) -- cycle;}
  \draw[bkLine, line width=0.4pt] (0,0) grid (4,4);
  \draw[bkBlue!80!black, line width=1.2pt] (0,0) -- (4,4);
  \draw[bkInk, line width=0.6pt] (0,0) -- (4.3,0) (0,0) -- (0,4.3);
  \foreach \k in {1,2,3,4} \node[font=\sffamily\scriptsize, anchor=north] at (\k,-0.1) {\k};
\end{tikzpicture}
\end{minipage}

\medskip
4〜8秒の長方形は、たて4マス、横4マスで $4\times4=16$ マスです。8〜10秒の台形は、まるごとのマスが $3+2=5$ こと、半分のマスが2こ（＝1マス）で6マスです。

全部で $8+16+6=30$ マスです。1マスが\qty{1}{m}なので、自転車は10秒で\qty{30}{m}進んだことになります。グラフの面積（マスの数）が、進んだ距離になっています。なぜそうなるのかを、この章で確かめていきます。

\begin{omoidasou}
\bkrecall{2-3}{道のり＝速さ×時間。速さが一定なら、$v$-$t$ グラフ（速さと時間のグラフ）の長方形の面積が道のり。}
\bkrecall{4-5}{三角形の面積＝底辺×高さ÷2（同じ三角形を2つ合わせると長方形になるから）。台形の面積＝（上底＋下底）×高さ÷2（同じ台形を2つ、さかさに合わせると平行四辺形になるから）。}
\bkrecall{5-5}{加速度は「1秒ごとに速度が何 m/s 変わるか」。加速度が一定の $a$ なら $v=v_0+at$、$x=x_0+v_0t+\dfrac{1}{2}at^2$。}
\bkrecall{5-8}{短い時間 $\Delta t$ のあいだは、速度はほとんど変わらない。だから $\Delta x\approx v\,\Delta t$（進む距離$\approx$速度×時間）。}
\bkrecall{2-10}{換算は「1をかける」。\qty{1000}{m/km}（1 kmあたり1000 m）、\qty{3600}{s/h}（1時間あたり3600秒）、\qty{60}{s/min}（1分あたり60秒）。}
\end{omoidasou}

%% ---------------- 数学の道具 ----------------
\bksection{doguu}{数学の道具}

\bksubsection{1}{速さが一定なら、長方形の面積}

\label{bk:5-9:chouhou}車が、\qty{6}{m/s}（1秒に\qty{6}{m}進む速さ）で2秒走ります。1秒めに\qty{6}{m}、2秒めにも\qty{6}{m}進むので、合わせて $6+6=6\times2=\qty{12}{m}$ です。

\begin{bkfigure}
\begin{tikzpicture}[x=1mm, y=1mm]
  % 数直線（1 m を 7 mm）
  \draw[bk axis] (-3,0) -- (92,0);
  \node[font=\sffamily\small, anchor=west] at (93,0) {m};
  \foreach \p in {0,6,12} {\draw[bkInk, line width=0.8pt] ({7*\p},-1.5) -- ({7*\p},1.5);
    \node[font=\sffamily\small, anchor=north] at ({7*\p},-2.5) {\p};}
  \foreach \p in {1,...,11} \draw[bkInk, line width=0.4pt] ({7*\p},-0.8) -- ({7*\p},0.8);
  \foreach \p/\t in {0/0,6/1,12/2} {\bkcar{({7*\p-3.6},3.5)}{0.3}{bkRoute2}
    \node[font=\sffamily\footnotesize, anchor=south] at ({7*\p+0.9},10) {\t 秒};}
  \draw[bk arrow, bkOrange] (3,-8) -- (39,-8);
  \draw[bk arrow, bkOrange] (45,-8) -- (81,-8);
  \node[font=\sffamily\footnotesize, anchor=north] at (21,-9.5) {1秒で\qty{6}{m}};
  \node[font=\sffamily\footnotesize, anchor=north] at (63,-9.5) {1秒で\qty{6}{m}};
\end{tikzpicture}
\bkfigcaption{\qty{6}{m/s}で2秒。1秒ごとに\qty{6}{m}ずつ進む。}
\end{bkfigure}

\noindent\begin{minipage}[c]{0.46\linewidth}\centering
\begin{tikzpicture}
\begin{axis}[bkplot, width=58mm, height=42mm, xmin=0, xmax=7.6, ymin=0, ymax=7.4, xtick={1,...,7}, ytick={3,6}, xlabel={$t$（s）}, ylabel={$v$（m/s）}]
  \fill[bkBlue!18] (axis cs:0,0) rectangle (axis cs:2,6);
  \fill[bkOrange!25] (axis cs:2,0) rectangle (axis cs:7,3);
  \addplot[bkcurve, line width=1.4pt] coordinates {(0,6) (2,6)};
  \addplot[bkcurve, line width=1.4pt] coordinates {(2,3) (7,3)};
  \node[font=\sffamily\small, fill=bkBlue!18, inner sep=1pt] at (axis cs:1,3.6) {\qty{12}{m}};
  \node[font=\sffamily\small, fill=bkOrange!25, inner sep=1pt] at (axis cs:4.5,1.5) {\qty{15}{m}};
\end{axis}
\end{tikzpicture}
\end{minipage}\hfill
\begin{minipage}[c]{0.5\linewidth}\raggedright
$v$-$t$ グラフでは、たて6（速さ）、横2（時間）の長方形の面積 $6\times2=12$ が、進んだ距離です。長方形を1秒ずつのたての列に切ると、1列が「1秒で進む距離」になっています。

単位も計算してみましょう。秒（s）どうしが約分されて、距離の単位 m が残ります。
\[ \qty{6}{m/s}\times\qty{2}{s}=6\,\frac{\mathrm{m}}{\cancel{\mathrm{s}}}\times2\,\cancel{\mathrm{s}}=\qty{12}{m} \]
\end{minipage}

\medskip
左のグラフのように、はじめの2秒は\qty{6}{m/s}、次の5秒は\qty{3}{m/s}で進んだとき、距離は
\[ 6\times2+3\times5=12+15=\qty{27}{m} \]
です。速さがとちゅうで変わっても、速さが一定の部分ごとに長方形に分けて、たし合わせればよいのです。

\bksubsection{2}{速さが変わり続けるとき：細かく分けて考える}

\label{bk:5-9:wakeru}止まっていた車が走り出して、1秒ごとに\qty{2}{m/s}ずつ速くなります（加速度\qty{2}{m/s^2}、5-5）。$t$ 秒後の速さは $v=2t$（m/s）です。1秒後は2、2秒後は4、4秒後は\qty{8}{m/s}です。

この車は、4秒で何m進むでしょう。

「速さ×時間」を使いたくても、速さがずっと変わっているので、どの速さをかければよいかわかりません。0をかければ0、8をかければ32。どちらもおかしいですね。

そこで、5-8の考え方を使います。\textbf{短い時間のあいだなら、速さはほとんど変わらない}ので、そのあいだは「速さ×時間」で距離を求めてよいのです。次の3つの段で計算します。

\begin{bkfigure}
\begin{tikzpicture}[x=1mm, y=1mm]
  \node[bk node, text width=30mm] (s1) at (0,0) {\textbf{① 分ける}\\4秒を1秒ずつ\\4つの区間に分ける};
  \node[bk node, text width=30mm] (s2) at (45,0) {\textbf{② かける}\\区間ごとに\\速さ×1秒};
  \node[bk node, text width=30mm] (s3) at (90,0) {\textbf{③ たす}\\4つの距離を\\全部たす};
  \draw[bk link] (s1) -- (s2);
  \draw[bk link] (s2) -- (s3);
\end{tikzpicture}
\bkfigcaption{速さが変わるときの距離の求め方。この3つの段を「積分」の考え方という。}
\end{bkfigure}

②で、区間のあいだの速さには、ひとまず\textbf{区間のはじめの速さ}を使います。たとえば1〜2秒の区間は、はじめ（1秒）の速さ\qty{2}{m/s}で1秒進んだと考えて、$2\times1=\qty{2}{m}$ です。

\begin{bkcenter}
\begin{bktable}{colspec={Q[l,m,wd=34mm] *{4}{Q[c,m,wd=14mm]} Q[c,m,wd=14mm]}, rowsep=1.5pt}
区間（s） & 0〜1 & 1〜2 & 2〜3 & 3〜4 & 合計 \\
はじめの速さ（m/s） & 0 & 2 & 4 & 6 & \\
速さ×1秒（m） & 0 & 2 & 4 & 6 & {\bfseries 12} \\
\end{bktable}
\bktabcaption{区間のはじめの速さを使うと、合計は\qty{12}{m}。}
\end{bkcenter}

\noindent 1つ1つの「速さ×1秒」は、$v$-$t$ グラフでは、はば1、高さが速さの\textbf{細い長方形}の面積です。それをならべると、階段の形になります（下の図の左）。0〜1秒の長方形は、高さが0なので見えません。

\begin{bkfigure}
\begin{tikzpicture}
\begin{axis}[bkplot, name=w1, width=60mm, height=54mm, xmin=0, xmax=4.7, ymin=0, ymax=9.4, xtick={1,2,3,4}, ytick={2,4,6,8}, xlabel={$t$}, ylabel={$v$}, title={\sffamily\small はじめの速さ（小さめ）}]
  \foreach \k in {1,2,3} {\edef\tmp{\noexpand\fill[bkBlue!25, draw=bkBlue!70, line width=0.5pt] (axis cs:\k,0) rectangle (axis cs:\k+1,{2*\k});}\tmp}
  \foreach \k in {0,1,2,3} {\edef\tmp{\noexpand\fill[bkOrange!50] (axis cs:\k,{2*\k}) -- (axis cs:\k+1,{2*\k}) -- (axis cs:\k+1,{2*\k+2}) -- cycle;}\tmp}
  \addplot[bkcurve, domain=0:4.2] {2*x};
  \foreach \k in {1,2,3} {\edef\tmp{\noexpand\node[font=\noexpand\sffamily\noexpand\footnotesize] at (axis cs:\k+0.5,\k) {\the\numexpr2*\k\relax};}\tmp}
\end{axis}
\begin{axis}[bkplot, at={(w1.east)}, anchor=west, xshift=12mm, width=60mm, height=54mm, xmin=0, xmax=4.7, ymin=0, ymax=9.4, xtick={1,2,3,4}, ytick={2,4,6,8}, xlabel={$t$}, ylabel={$v$}, title={\sffamily\small おわりの速さ（大きめ）}]
  \foreach \k in {0,1,2,3} {\edef\tmp{\noexpand\fill[bkBlue!25, draw=bkBlue!70, line width=0.5pt] (axis cs:\k,0) rectangle (axis cs:\k+1,{2*\k+2});}\tmp}
  \foreach \k in {0,1,2,3} {\edef\tmp{\noexpand\fill[bkRed!35] (axis cs:\k,{2*\k}) -- (axis cs:\k,{2*\k+2}) -- (axis cs:\k+1,{2*\k+2}) -- cycle;}\tmp}
  \addplot[bkcurve, domain=0:4.2] {2*x};
  \node[font=\sffamily\footnotesize] at (axis cs:0.72,0.5) {2};
  \foreach \k in {1,2,3} {\edef\tmp{\noexpand\node[font=\noexpand\sffamily\noexpand\footnotesize] at (axis cs:\k+0.5,\k) {\the\numexpr2*\k+2\relax};}\tmp}
\end{axis}
\end{tikzpicture}
\bkfigcaption{左：はじめの速さの階段（合計12）。オレンジは、数えもれたすき間。\\右：おわりの速さの階段（合計20）。赤は、はみ出した分。}
\end{bkfigure}

\noindent\textbf{12 mは少なすぎる}\quad 車は、どの区間でも、とちゅうで速くなっています。たとえば1〜2秒のあいだ、速さは2より大きく、4より小さい値です。それなのに、ずっと2のままとして計算したので、ほんとうより少なくなります。グラフでは、階段と直線のあいだのオレンジのすき間が、数えもれた分です。

\noindent\textbf{おわりの速さなら多すぎる}\quad 反対に、区間の\textbf{おわりの速さ}（1秒で2、2秒で4、3秒で6、4秒で8）を使うと、合計は
\[ 2+4+6+8=\qty{20}{m} \]
です。こんどはずっと速いとして計算したので、ほんとうより多くなります（右の図の赤い部分がはみ出した分）。

だから、ほんとうの距離は、12 mより多く、20 mより少ないことがわかります。
\[ 12<\text{ほんとうの距離}<20 \]

\bksubsection{3}{区間を短くすると、階段は三角形に近づく}

\label{bk:5-9:chikazuku}区間をもっと短くしてみましょう。区間のはばを $\Delta t$（デルタ・ティー）と書きます。$\Delta t=0.5$ なら、4秒は8つの区間に分かれます。1つめの区間のはじめの速さは0、2つめは $2\times0.5=1$、3つめは $2\times1=2$……です。

\begin{align*}
\text{はじめの速さ：}\quad&(0+1+2+3+4+5+6+7)\times0.5=28\times0.5=14\\
\text{おわりの速さ：}\quad&(1+2+3+4+5+6+7+8)\times0.5=36\times0.5=18
\end{align*}
0.5は、どの区間にも共通の「×0.5秒」なので、先に速さをたしてから、まとめて0.5をかけました（分配法則、2-11）。同じように $\Delta t$ を小さくして計算すると、次の表になります。

\begin{bkcenter}
\begin{bktable}{colspec={Q[c,m,wd=17mm] Q[c,m,wd=14mm] Q[c,m,wd=30mm] Q[c,m,wd=30mm] Q[c,m,wd=16mm]}, rowsep=1.5pt}
$\Delta t$（s） & 区間の数 & はじめの速さの和（m） & おわりの速さの和（m） & 2つの差 \\
1 & 4 & 12 & 20 & 8 \\
0.5 & 8 & 14 & 18 & 4 \\
0.25 & 16 & 15 & 17 & 2 \\
0.1 & 40 & 15.6 & 16.4 & 0.8 \\
0.01 & 400 & 15.96 & 16.04 & 0.08 \\
\end{bktable}
\bktabcaption{区間を短くすると、2つの和はどちらも16に近づく。}
\end{bkcenter}

$\Delta t$ を小さくするほど、小さめの和は大きくなり、大きめの和は小さくなって、2つの差は0に近づきます。2つの和は、どちらも16に近づいていきます。ほんとうの距離は、いつもこの2つのあいだにあるので、\textbf{ほんとうの距離は\qty{16}{m}}です。

図でも見てみましょう。はじめの速さの階段で、区間を短くしていきます。

\begin{bkfigure}
\begin{tikzpicture}[x=1mm, y=1mm]
  \foreach \n/\px/\dtv/\sm/\gp in {4/0/1/12/4, 8/46/0.5/14/2, 16/92/0.25/15/1} {
    \begin{scope}[shift={(\px,0)}]
      \foreach \k in {1,...,\n} {
        \fill[bkBlue!25] ({32*(\k-1)/\n},0) rectangle ({32*\k/\n},{24*(\k-1)/\n});
        \fill[bkOrange!55] ({32*(\k-1)/\n},{24*(\k-1)/\n}) -- ({32*\k/\n},{24*(\k-1)/\n}) -- ({32*\k/\n},{24*\k/\n}) -- cycle;}
      \foreach \k in {2,...,\n} \draw[bkBlue!70, line width=0.3pt] ({32*(\k-1)/\n},0) rectangle ({32*\k/\n},{24*(\k-1)/\n});
      \draw[bkBlue!80!black, line width=1pt] (0,0) -- (33,24.75);
      \draw[bkInk, line width=0.6pt, -{Stealth[length=3pt]}] (0,0) -- (37,0);
      \draw[bkInk, line width=0.6pt, -{Stealth[length=3pt]}] (0,0) -- (0,28);
      \node[font=\sffamily\scriptsize, anchor=north, inner sep=1pt] at (32,-0.5) {4};
      \node[font=\sffamily\bfseries\footnotesize, anchor=north west, inner sep=0pt] at (0,-5) {$\Delta t=\dtv$};
      \node[font=\sffamily\footnotesize, anchor=north west, inner sep=0pt, align=left] at (0,-10.5) {階段 \sm\ m\\{\color{bkOrange!80!black}すき間 \gp\ m}};
    \end{scope}}
\end{tikzpicture}
\bkfigcaption{区間を短くすると、オレンジのすき間はどんどん小さくなり、階段は三角形に近づく。}
\end{bkfigure}

\noindent\textbf{すき間が0に近づくわけ}\quad オレンジのすき間は、小さな三角形の集まりです。1つの三角形は、底辺が $\Delta t$、高さが「$\Delta t$ のあいだに速さがふえた分」の $2\times\Delta t$ です。
\begin{itemize}[leftmargin=1.6\zw, itemsep=1pt, topsep=2pt, label={\color{bkTheme}●}]
\item 小さな三角形1つの面積は $\Delta t\times2\Delta t\div2=\Delta t\times\Delta t$
\item 三角形の数は、区間の数と同じ $4\div\Delta t$ こ
\item すき間の合計は $(4\div\Delta t)\times\Delta t\times\Delta t=4\times\Delta t$
\end{itemize}
$\Delta t=1$ なら4、$\Delta t=0.5$ なら2、$\Delta t=0.25$ なら1で、図と合っています。$\Delta t$ を半分にすると、すき間の合計も半分になります。$\Delta t$ を0に近づけると、すき間は0に近づき、階段の面積は、直線の下の\textbf{三角形の面積}にかぎりなく近づきます。

\begin{taisetsu}
① 分ける\quad ② かける（速さ×$\Delta t$）\quad ③ たす\\
$\Delta t$ を小さくしていくと、たした値は、$v$-$t$ グラフの下の面積に近づく。\\
その面積が、ほんとうに進んだ距離。
\end{taisetsu}

\begin{column}{表計算で階段の和を計算しよう}
②・③の車（$v=2t$）を、表計算で $\Delta t=0.1$ にして計算します（2-18）。A列に区間のはじめの時刻、B列にその速さ、C列に「速さ×0.1」、D列にC列の合計をつみ上げます。
\begin{bkcenter}
\bksheet{4}{22}{%
1 & はじめの時刻 & 速さ $2t$ & 速さ×0.1 & 合計 \\
2 & 0 & 0 & 0 & 0 \\
3 & 0.1 & 0.2 & 0.02 & 0.02 \\
4 & 0.2 & 0.4 & 0.04 & 0.06 \\
5 & （下へコピー） & & & \\
41 & 3.9 & 7.8 & 0.78 & 15.6 \\}
\end{bkcenter}
B2に \texttt{=2*A2}、C2に \texttt{=B2*0.1}、D2に \texttt{=C2}、D3に \texttt{=D2+C3} と入れて下へコピーすると、41行めの合計は15.6になり、③の表と合います。0.1を0.01にして、行を400行にすると15.96。ほんとうの16に近づいていきます。
\end{column}

\Needspace{6\baselineskip}
\bksubsection{4}{一定の割合で速くなるなら、三角形の面積}

\label{bk:5-9:sankaku}②・③の車の三角形は、底辺4（時間）、高さ8（4秒後の速さ）なので、
\[ 4\times8\div2=\qty{16}{m} \]
表で近づいていった16と、ぴったり同じです。

\noindent\begin{minipage}[c]{0.52\linewidth}\raggedright
\textbf{なぜ÷2なの？}\quad 右の図のように、同じ三角形を2つ合わせると、たて8、横4の長方形になります。長方形の面積 $4\times8=32$ は「ずっと\qty{8}{m/s}で走った距離」で、三角形はその半分です。
\end{minipage}\hfill
\begin{minipage}[c]{0.44\linewidth}\centering
\begin{tikzpicture}
\begin{axis}[bkplot, width=56mm, height=50mm, xmin=0, xmax=4.9, ymin=0, ymax=9.6, xtick={2,4}, ytick={4,8}, xlabel={$t$}, ylabel={$v$}]
  \fill[bkGray!20] (axis cs:0,0) -- (axis cs:0,8) -- (axis cs:4,8) -- cycle;
  \fill[bkBlue!28] (axis cs:0,0) -- (axis cs:4,8) -- (axis cs:4,0) -- cycle;
  \draw[bkGray, line width=0.6pt, dashed] (axis cs:0,8) -- (axis cs:4,8) -- (axis cs:4,0);
  \addplot[bkcurve, domain=0:4.3] {2*x};
  \node[font=\sffamily\footnotesize, fill=bkBlue!28, inner sep=1pt] at (axis cs:2.85,1.6) {\qty{16}{m}};
  \node[font=\sffamily\footnotesize, text=bkGray!80!black] at (axis cs:1.15,6.4) {同じ形};
\end{axis}
\end{tikzpicture}
\end{minipage}

\medskip
\noindent\begin{minipage}[c]{0.44\linewidth}\centering
\begin{tikzpicture}
\begin{axis}[bkplot, width=56mm, height=50mm, xmin=0, xmax=4.9, ymin=0, ymax=9.6, xtick={2,4}, ytick={4,8}, xlabel={$t$}, ylabel={$v$}]
  \fill[bkBlue!28] (axis cs:0,0) rectangle (axis cs:4,4);
  \fill[bkOrange!55] (axis cs:2,4) -- (axis cs:4,4) -- (axis cs:4,8) -- cycle;
  \fill[white] (axis cs:0,0) -- (axis cs:2,4) -- (axis cs:0,4) -- cycle;
  \fill[bkOrange!25] (axis cs:0,0) -- (axis cs:2,4) -- (axis cs:0,4) -- cycle;
  \draw[bkGray, line width=0.6pt, dashed] (axis cs:0,4) -- (axis cs:4,4);
  \addplot[bkcurve, domain=0:4.3] {2*x};
  \node[font=\sffamily\scriptsize, text=bkOrange!80!black] at (axis cs:3.35,5.5) {A};
  \node[font=\sffamily\scriptsize, text=bkOrange!80!black] at (axis cs:0.6,2.85) {B};
\end{axis}
\end{tikzpicture}
\end{minipage}\hfill
\begin{minipage}[c]{0.52\linewidth}\raggedright
\textbf{平均の速さで考える}\quad 速さは0から8まで、まっすぐふえていきます。2秒の速さ4をまん中にして、1秒の速さ2と3秒の速さ6、0.5秒の速さ1と3.5秒の速さ7……と、前と後ろで組にすると、どの組も平均は4です。だから、4秒間の\textbf{平均の速さは $(0+8)\div2=\qty{4}{m/s}$}。
\[ 4\times4=\qty{16}{m} \]
図では、三角形の上の部分Aを切りとって、左上のBにうめると、たて4、横4の長方形になります。2通りの考え方で、同じ16 mになりました。
\end{minipage}

\medskip
\noindent\textbf{文字の式にする}\quad 止まっていた物が、一定の加速度 $a$ で $t$ 秒間速くなるとします。
\begin{itemize}[leftmargin=1.6\zw, itemsep=1pt, topsep=2pt, label={\color{bkTheme}●}]
\item 三角形の底辺は、時間の $t$。
\item 高さは、$t$ 秒後の速さ。1秒に $a$ ずつ、$t$ 秒ふえるので $at$。
\item ÷2 は、三角形が長方形の半分だから。$\div2$ は $\times\dfrac{1}{2}$ と同じ。
\end{itemize}
\[ x=t\times at\times\frac{1}{2}=\frac{1}{2}at^2 \]
これは、5-5で微分して確かめた式 $x=\dfrac{1}{2}at^2$ と同じです。$a=2$、$t=4$ を入れると $\dfrac{1}{2}\times2\times16=16$ で、②・③の車とも合います。

\Needspace{12\baselineskip}
\bksubsection{5}{はじめから動いていれば、台形の面積}

\label{bk:5-9:daikei}\qty{5}{m/s}で走っていた自転車が、一定の割合で速くなり、4秒後に\qty{15}{m/s}になりました。4秒で何m進むでしょう。$v$-$t$ グラフは、たての辺が5と15、はばが4の\textbf{台形}になります。

\noindent\begin{minipage}[c]{0.46\linewidth}\centering
\begin{tikzpicture}
\begin{axis}[bkplot, width=58mm, height=54mm, xmin=0, xmax=5.2, ymin=0, ymax=18, xtick={4}, ytick={5,15}, xlabel={$t$}, ylabel={$v$}]
  \fill[bkBlue!25] (axis cs:0,0) rectangle (axis cs:4,5);
  \fill[bkOrange!40] (axis cs:0,5) -- (axis cs:4,5) -- (axis cs:4,15) -- cycle;
  \draw[bkGray, line width=0.6pt, dashed] (axis cs:4,0) -- (axis cs:4,15);
  \addplot[bkcurve, domain=0:4.4] {5+2.5*x};
  \node[font=\sffamily\footnotesize] at (axis cs:2,2.5) {\qty{20}{m}};
  \node[font=\sffamily\footnotesize] at (axis cs:3.05,8) {\qty{20}{m}};
\end{axis}
\end{tikzpicture}
\end{minipage}\hfill
\begin{minipage}[c]{0.5\linewidth}\raggedright
\textbf{分けて考える}\quad 台形を、下の長方形（青）と上の三角形（オレンジ）に分けます。
\begin{itemize}[leftmargin=1.2\zw, itemsep=1pt, topsep=2pt, label={\color{bkTheme}●}]
\item 長方形：はじめの速さ5のまま4秒走った分。$5\times4=20$
\item 三角形：速くなったことで、よけいに進んだ分。高さは $15-5=10$（ふえた速さ）。$4\times10\div2=20$
\end{itemize}
合わせて $20+20=\qty{40}{m}$ です。
\end{minipage}

\medskip
\noindent\textbf{台形の公式で}\quad このグラフの台形は、横にねています。たての2つの辺（5と15）が上底と下底、横のはば（4秒）が高さです。
\[ (5+15)\times4\div2=20\times4\div2=\qty{40}{m} \]
\textbf{平均の速さで}\quad 速さは5から15までまっすぐふえるので、平均の速さは $(5+15)\div2=\qty{10}{m/s}$。$10\times4=\qty{40}{m}$。3通りとも同じです。

\noindent\textbf{文字の式にする}\quad はじめの速度 $v_0$、加速度 $a$ で $t$ 秒走ると、$t$ 秒後の速度は $v_0+at$ です。長方形は $v_0\times t$、三角形は高さ $at$ なので $\dfrac{1}{2}at^2$（④と同じ）。
\[ x=v_0t+\frac{1}{2}at^2 \]
5-5の $x=x_0+v_0t+\dfrac{1}{2}at^2$ の、$x_0$（はじめの位置）以外の部分です。面積が表すのは、はじめの位置から\textbf{どれだけ変わったか}なので、$x_0$ は出てきません。

\bksubsection{6}{おそくなるときも、同じ}

\label{bk:5-9:osoku}ブレーキをかけて速さがへるとき、$v$-$t$ グラフは右下がりの直線になります。形は、向きがちがうだけの三角形や台形です。

\begin{bkfigure}
\begin{tikzpicture}
\begin{axis}[bkplot, name=o1, width=58mm, height=46mm, xmin=0, xmax=6.2, ymin=0, ymax=24, xtick={5}, ytick={20}, xlabel={$t$}, ylabel={$v$}, title={\sffamily\small 止まるまで（三角形）}]
  \fill[bkBlue!25] (axis cs:0,0) -- (axis cs:0,20) -- (axis cs:5,0) -- cycle;
  \addplot[bkcurve, domain=0:5] {20-4*x};
  \node[font=\sffamily\footnotesize] at (axis cs:1.5,5) {\qty{50}{m}};
\end{axis}
\begin{axis}[bkplot, at={(o1.east)}, anchor=west, xshift=12mm, width=58mm, height=46mm, xmin=0, xmax=5.2, ymin=0, ymax=18, xtick={4}, ytick={5,15}, xlabel={$t$}, ylabel={$v$}, title={\sffamily\small とちゅうまで（台形）}]
  \fill[bkBlue!25] (axis cs:0,0) -- (axis cs:0,15) -- (axis cs:4,5) -- (axis cs:4,0) -- cycle;
  \draw[bkGray, line width=0.6pt, dashed] (axis cs:4,0) -- (axis cs:4,5);
  \addplot[bkcurve, domain=0:4] {15-2.5*x};
  \node[font=\sffamily\footnotesize] at (axis cs:2,3.6) {\qty{40}{m}};
\end{axis}
\end{tikzpicture}
\bkfigcaption{左：\qty{20}{m/s}から5秒で止まる。右：\qty{15}{m/s}から4秒で\qty{5}{m/s}になる。}
\end{bkfigure}

\begin{itemize}[leftmargin=1.6\zw, itemsep=1pt, topsep=2pt, label={\color{bkTheme}●}]
\item 左：底辺5、高さ20の三角形なので、$5\times20\div2=\qty{50}{m}$。\\平均の速さで考えると、$(20+0)\div2=10$、$10\times5=50$。\\5-5の例（\qty{20}{m/s}の車が加速度 $\qty{-4}{m/s^2}$ で止まるまで）と同じ結果です。
\item 右：上底15、下底5、高さ4の台形。$(15+5)\times4\div2=\qty{40}{m}$。⑤の速くなる自転車（5から15）と、ちょうど左右を入れかえた形なので、面積も同じ40です。
\end{itemize}

\Needspace{6\baselineskip}
\bksubsection{7}{速度がマイナスの部分は、マイナスの面積}

\label{bk:5-9:minus}ろう下で、右向きを＋とします。\qty{4}{m/s}で3秒進み、そのあと\qty{-2}{m/s}で2秒もどりました。\qty{-2}{m/s}は「－の向き（左）に、1秒に\qty{2}{m}」という意味です（5-5）。

\noindent\begin{minipage}[c]{0.5\linewidth}\centering
\begin{tikzpicture}
% 第7回：目もりの数字 3・5 は、赤い四角の中ではなく下に出し、点線でつなぐ
\begin{axis}[bkplot, width=62mm, height=46mm, xmin=0, xmax=6, ymin=-3.6, ymax=5, xtick={3,5}, xticklabels={}, ytick={-2,4}, xlabel={$t$（s）}, ylabel={$v$（m/s）}]
  \fill[bkBlue!20] (axis cs:0,0) rectangle (axis cs:3,4);
  \fill[bkRed!25] (axis cs:3,0) rectangle (axis cs:5,-2);
  \addplot[bkcurve, line width=1.4pt] coordinates {(0,4) (3,4)};
  \addplot[bkcurve, line width=1.4pt] coordinates {(3,-2) (5,-2)};
  \draw[bkGray, line width=0.5pt, dash pattern=on 1.5pt off 1pt] (axis cs:3,-2) -- (axis cs:3,-2.55) (axis cs:5,-2) -- (axis cs:5,-2.55);
  \node[font=\footnotesize, text=bkInk, anchor=north] at (axis cs:3,-2.6) {$3$};
  \node[font=\footnotesize, text=bkInk, anchor=north] at (axis cs:5,-2.6) {$5$};
  \node[font=\sffamily\small] at (axis cs:1.5,2) {$+12$};
  \node[font=\sffamily\small] at (axis cs:4,-1) {$-4$};
\end{axis}
\end{tikzpicture}
\end{minipage}\hfill
\begin{minipage}[c]{0.46\linewidth}\raggedright
横軸より上の長方形は、たて4、横3なので、面積は $4\times3=12$。右へ\qty{12}{m}進んだ、という意味です。

横軸より下の長方形は、たて $-2$、横2なので、面積を $(-2)\times2=-4$ と考えます。左へ\qty{4}{m}もどった、という意味です。
\end{minipage}

\begin{bkfigure}
\begin{tikzpicture}[x=1mm, y=1mm]
  \draw[bk axis] (-4,0) -- (104,0);
  \node[font=\sffamily\small, anchor=west] at (105,0) {m};
  \foreach \p in {0,...,13} \draw[bkInk, line width=0.4pt] ({7*\p},-0.9) -- ({7*\p},0.9);
  \foreach \p in {0,4,8,12} \node[font=\sffamily\small, anchor=north] at ({7*\p},-2) {\p};
  \fill[bkInk] (0,0) circle (1.3);
  \fill[bkRed] (56,0) circle (1.3);
  \draw[bk arrow, bkBlue] (0,5) -- (84,5);
  \node[font=\sffamily\footnotesize, anchor=south] at (36,6.5) {3秒で $+\qty{12}{m}$};
  \draw[bk arrow, bkRed] (84,13) -- (56,13);
  \node[font=\sffamily\footnotesize, anchor=south] at (70,14.5) {2秒で $-\qty{4}{m}$};
  \node[font=\sffamily\footnotesize, anchor=north] at (0,-7) {はじめ};
  \node[font=\sffamily\footnotesize, anchor=north, text=bkRed!80!black] at (56,-7) {最後の位置};
\end{tikzpicture}
\bkfigcaption{はじめの位置0から、12まで進み、8までもどった。}
\end{bkfigure}

最後の位置は、はじめの位置から $12+(-4)=\qty{8}{m}$ 右です。このように、はじめの位置から最後の位置までの変化を\textbf{変位}（へんい）といいます。変位は、向きをふくめた量で、上の面積の合計（マイナスの面積もそのままたす）が変位です。\bkidx{へんい}{変位}

一方、実際に動いた長さは、もどった分もたして $12+4=\qty{16}{m}$ です。これを\textbf{道のり}といいます。道のりは、マイナスの面積も、プラスの大きさにしてたします。\bkidx{みちのり}{道のり}

\noindent\textbf{線が横軸を横切るとき}\quad ボールを\qty{20}{m/s}で真上に投げ上げます（上向きを＋）。速度は $v=20-10t$（5-5）。2秒までは＋（上がる）、2秒で0（いちばん高い）、2秒からは－（落ちる）です。

\noindent\begin{minipage}[c]{0.42\linewidth}\centering
\begin{tikzpicture}
\begin{axis}[bkplot, width=52mm, height=50mm, xmin=0, xmax=4.8, ymin=-24, ymax=26, xtick={1}, ytick={-20,20}, xlabel={$t$}, ylabel={$v$}]
  \fill[bkBlue!25] (axis cs:0,0) -- (axis cs:0,20) -- (axis cs:2,0) -- cycle;
  \fill[bkRed!25] (axis cs:2,0) -- (axis cs:4,-20) -- (axis cs:4,0) -- cycle;
  \draw[bkGray, line width=0.4pt] (axis cs:3,0) -- (axis cs:3,1.2) (axis cs:4,0) -- (axis cs:4,1.2);
  \addplot[bkcurve, domain=0:4.2] {20-10*x};
  \node[font=\sffamily\footnotesize] at (axis cs:0.6,5) {$+20$};
  \node[font=\sffamily\footnotesize] at (axis cs:3.55,-6.5) {$-20$};
  \node[font=\sffamily\footnotesize, anchor=south] at (axis cs:2,1.2) {2};
  \node[font=\sffamily\footnotesize, anchor=south] at (axis cs:3,1.2) {3};
  \node[font=\sffamily\footnotesize, anchor=south] at (axis cs:4,1.2) {4};
\end{axis}
\end{tikzpicture}
\end{minipage}\hfill
\begin{minipage}[c]{0.54\linewidth}\raggedright
\begin{itemize}[leftmargin=1.2\zw, itemsep=1pt, topsep=2pt, label={\color{bkTheme}●}]
\item 0〜2秒は、横軸より上の三角形。\\$2\times20\div2=20$ で、\qty{20}{m}上がる。
\item 2〜4秒は、横軸より下の三角形。\\$-(2\times20\div2)=-20$ で、\qty{20}{m}落ちる。
\item 4秒までの変位は $20+(-20)=0$。\\投げた所にもどってきた。
\item 道のりは $20+20=\qty{40}{m}$。
\end{itemize}
\end{minipage}

\Needspace{12\baselineskip}
\bksubsection{8}{面積をつみ上げて、位置のグラフをつくる}

\label{bk:5-9:tsumiage}面積は、とちゅうまでの分も計算できます。やってみようの自転車で、2秒ごとの面積を求め、0秒からの合計（出発点からの位置）を表にします。

\begin{bkcenter}
\begin{bktable}{colspec={Q[l,m,wd=30mm] *{6}{Q[c,m,wd=12mm]}}, rowsep=1.5pt, colsep=4pt, row{3}={font=\small}}
時刻 $t$（s） & 0 & 2 & 4 & 6 & 8 & 10 \\
速さ $v$（m/s） & 0 & 2 & 4 & 4 & 4 & 2 \\
2秒ごとの形 & & 三角形 & 台形 & 長方形 & 長方形 & 台形 \\
2秒ごとの面積（m） & & 2 & 6 & 8 & 8 & 6 \\
0秒からの合計＝位置 $x$（m） & 0 & 2 & 8 & 16 & 24 & 30 \\
\end{bktable}
\bktabcaption{2秒ごとの面積をつみ上げると、位置になる。}
\end{bkcenter}

\noindent 2秒ごとの面積の求め方：0〜2秒は三角形 $2\times2\div2=2$。2〜4秒は台形 $(2+4)\times2\div2=6$。4〜6秒と6〜8秒は長方形 $4\times2=8$。8〜10秒は台形 $(4+2)\times2\div2=6$。合計の行は、前の合計に次の面積をたしていきます（$0+2=2$、$2+6=8$、$8+8=16$……）。

\noindent\begin{minipage}[c]{0.52\linewidth}\centering
\begin{tikzpicture}
\begin{axis}[bkplot, width=64mm, height=46mm, xmin=0, xmax=11, ymin=0, ymax=34, xtick={2,4,6,8,10}, ytick={10,20,30}, xlabel={$t$（s）}, ylabel={$x$（m）}]
  \addplot[bkcurve, domain=0:4] {x^2/2};
  \addplot[bkcurve, domain=4:8] {8+4*(x-4)};
  \addplot[bkcurve, domain=8:10] {24+4*(x-8)-0.5*(x-8)^2};
  \addplot[only marks, mark=*, mark size=1.5pt, bkInk] coordinates {(0,0) (2,2) (4,8) (6,16) (8,24) (10,30)};
\end{axis}
\end{tikzpicture}
\end{minipage}\hfill
\begin{minipage}[c]{0.44\linewidth}\raggedright
表の点を $x$-$t$ グラフにすると、左のようになります。速さがふえている0〜4秒はだんだん急に、速さ一定の4〜8秒はまっすぐ、おそくなる8〜10秒はだんだんゆるやかになります。

4〜8秒の傾きは\\$(24-8)\div(8-4)=4$ で、このときの速さ\qty{4}{m/s}と同じです。
\end{minipage}

\medskip
$x$-$t$ グラフの傾きは、もとの速さになっています（5-3・5-5で学んだ「位置を微分すると速度」）。面積をつみ上げる計算と、傾きを求める計算は、ちょうど反対向きの関係になっています。

\begin{bkcenter}
\begin{bktable}{colspec={Q[l,m,wd=22mm] Q[l,m,wd=50mm] Q[l,m,wd=50mm]}, rowsep=2pt, cells={font=\small}, row{1}={font=\sffamily\bfseries\small}}
 & 微分（5-1〜5-8） & 面積をつみ上げる（5-9〜） \\
向き & 位置 $x$ → 速度 $v$ & 速度 $v$ → 位置の変化 \\
短い区間で & わり算 $\Delta x\div\Delta t$ & かけ算 $v\times\Delta t$ \\
そのあと & $\Delta t$ を0に近づける & たして、$\Delta t$ を0に近づける \\
グラフでは & 傾き & 面積 \\
\end{bktable}
\bktabcaption{微分と、面積をつみ上げる計算のくらべ。5-11で「積分は微分の逆」として整理する。}
\end{bkcenter}

\bksubsection{9}{はてな？の答え：タコグラフ}

\label{bk:5-9:taco}\bkhatenaref のグラフは、時間の単位が h（時間）、速さの単位が km/h です。km/h×h＝km なので、面積がそのまま km になります。単位の換算はいりません。グラフを、まっすぐな線ごとの区間に分けます。

\begin{bkcenter}
\begin{bktable}{colspec={Q[c,m,wd=20mm] Q[c,m,wd=16mm] Q[c,m,wd=48mm] Q[c,m,wd=16mm]}, rowsep=1.5pt}
区間（h） & 形 & 面積の計算 & 距離（km） \\
0〜0.5 & 三角形 & $0.5\times40\div2$ & 10 \\
0.5〜2 & 長方形 & $40\times1.5$ & 60 \\
2〜2.5 & 三角形 & $0.5\times40\div2$ & 10 \\
2.5〜3 & （休けい） & 速さ0 & 0 \\
3〜3.5 & 三角形 & $0.5\times60\div2$ & 15 \\
3.5〜5.5 & 長方形 & $60\times2$ & 120 \\
5.5〜6 & 三角形 & $0.5\times60\div2$ & 15 \\
\end{bktable}
\bktabcaption{全部で $10+60+10+0+15+120+15=230$。}
\end{bkcenter}

\begin{bkfigure}
\begin{tikzpicture}
\begin{axis}[bkplot, width=110mm, height=44mm, xmin=0, xmax=6.6, ymin=0, ymax=72, xtick={0.5,1,...,6}, xticklabels={0.5,1,1.5,2,2.5,3,3.5,4,4.5,5,5.5,6}, ytick={20,40,60}, xlabel={$t$（h）}, ylabel={$v$（km/h）}, x tick label style={font=\sffamily\scriptsize}]
  \fill[bkYellow!45] (axis cs:0,0) -- (axis cs:0.5,40) -- (axis cs:0.5,0) -- cycle;
  \fill[bkBlue!22] (axis cs:0.5,0) rectangle (axis cs:2,40);
  \fill[bkYellow!45] (axis cs:2,0) -- (axis cs:2,40) -- (axis cs:2.5,0) -- cycle;
  \fill[bkYellow!45] (axis cs:3,0) -- (axis cs:3.5,60) -- (axis cs:3.5,0) -- cycle;
  \fill[bkBlue!22] (axis cs:3.5,0) rectangle (axis cs:5.5,60);
  \fill[bkYellow!45] (axis cs:5.5,0) -- (axis cs:5.5,60) -- (axis cs:6,0) -- cycle;
  \addplot[bkcurve, line width=1.3pt] coordinates {(0,0) (0.5,40) (2,40) (2.5,0) (3,0) (3.5,60) (5.5,60) (6,0)};
  \node[font=\sffamily\footnotesize] at (axis cs:1.25,20) {60};
  \node[font=\sffamily\footnotesize] at (axis cs:4.5,30) {120};
  \node[font=\sffamily\scriptsize] at (axis cs:0.36,8) {10};
  \node[font=\sffamily\scriptsize] at (axis cs:2.14,8) {10};
  \node[font=\sffamily\scriptsize] at (axis cs:3.36,10) {15};
  \node[font=\sffamily\scriptsize] at (axis cs:5.64,10) {15};
\end{axis}
\end{tikzpicture}
\bkfigcaption{タコグラフの面積を、三角形（黄）と長方形（青）に分けた。数は面積（km）。}
\end{bkfigure}

\noindent このトラックは、6時間で\qty{230}{km}走りました。速さだけの記録から、走った距離がわかりました。ほんとうのタコグラフの線はもっとでこぼこですが、細かく分けて（5-10）たせば、同じように距離が求められます。

\begin{motto}[2]
\textbf{グラフが曲がっていても、マスを数えればおよそわかる}

\begin{bkfigure}
\begin{tikzpicture}[x=6.5mm, y=4.4mm]
  \foreach \k/\lo/\hi in {0/0/3,1/2/5,2/4/7,3/6/8,4/7/8,5/7/8,6/6/8,7/4/7,8/2/5,9/0/3} {
    \fill[bkYellow!55] (\k,0) rectangle (\k+1,\lo);
    \fill[bkOrange!35] (\k,\lo) rectangle (\k+1,\hi);}
  \draw[bkLine, line width=0.4pt] (0,0) grid (10,8);
  \draw[bkBlue!80!black, line width=1.3pt, domain=0:10, samples=60] plot (\x,{3*\x-0.3*\x*\x});
  \draw[bkInk, line width=0.6pt, -{Stealth[length=3pt]}] (0,0) -- (10.7,0);
  \draw[bkInk, line width=0.6pt, -{Stealth[length=3pt]}] (0,0) -- (0,8.7);
  \foreach \k in {2,4,6,8,10} \node[font=\sffamily\scriptsize, anchor=north] at (\k,-0.15) {\k};
  \foreach \k in {2,4,6,8} \node[font=\sffamily\scriptsize, anchor=east] at (-0.15,\k) {\k};
  \node[font=\sffamily\scriptsize, anchor=west] at (10.7,0) {$t$（s）};
  \node[font=\sffamily\scriptsize, anchor=south] at (0,8.7) {$v$（m/s）};
\end{tikzpicture}
\bkfigcaption{黄：まるごと入っているマス（38こ）。オレンジ：線がとおり、一部だけ入っているマス（24こ）。}
\end{bkfigure}

ほんとうの乗り物の $v$-$t$ グラフは、上の図のように曲がっています。それでも、マスを数えれば、およその距離がわかります。一部だけ入っているマスを、ならして半分と見ると、$38+24\div2=50$ で、約\qty{50}{m}です。

\medskip
この図の速さは $v=3t-0.3t^2$ で、5-10・5-11の方法で計算すると、ちょうど\qty{50}{m}になります。マスを細かくするほど、数えた値はほんとうの値に近づきます。「細かく分けて、たす」は、曲がったグラフでもそのまま使えるのです。
\end{motto}

\begin{column}{ガリレオより前の発見}
一定の割合で速くなるときに進む距離は、「はじめの速さとおわりの速さの平均で、ずっと走ったときの距離」に等しくなります。④・⑤の「平均の速さで考える」と同じことです。このきまりは、700年ほど前に、イギリスのオックスフォードの学者たちが見つけ、フランスのオレームが図（いまの $v$-$t$ グラフのようなもの）を使って説明しました。ガリレオは、これを実験で確かめ、落ちる物の運動にあてはめました。
\end{column}

\begin{taisetsu}
$v$-$t$ グラフと横軸のあいだの面積 ＝ そのあいだの位置の変化（変位）\\
速さが変わるときは、細かく分けて、かけて（速さ×$\Delta t$）、たす。\\
横軸より下の部分は、マイナスの面積として数える。
\end{taisetsu}

\begin{motto}[2]
\textbf{$a$-$t$ グラフの面積は、速度の変化}

加速度は「1秒あたりの速度の変化」でした（5-5）。だから、加速度のグラフ（$a$-$t$ グラフ）の面積は、そのあいだの速度の変化になります。単位も、$\unit{m/s^2}\times\unit{s}=\unit{m/s}$ と、速度の単位になります。加速度が\qty{2}{m/s^2}で3秒なら、$2\times3=6$ で、速度は\qty{6}{m/s}ふえます。スマートフォンの加速度センサーの値から速度を計算するのは、この考え方です。
\end{motto}

\begin{kotoba}
\small
\bkword{$v$-$t$ グラフ}{横軸に時間、たて軸に速度をとったグラフ。面積が位置の変化になる。}
\bkword{変位}{はじめの位置から、最後の位置までの変化。向きをふくむ（マイナスにもなる）。}
\bkword{道のり}{実際に動いた長さの合計。向きは考えない（いつもプラス）。}
\end{kotoba}

%% ---------------- 例題 ----------------
\bksection{reidai}{例題でたしかめよう}

\begin{reidai}[長方形に分ける]
電車が、\qty{20}{m/s}で30秒走ったあと、\qty{10}{m/s}で20秒走りました。進んだ距離を求めましょう。\par
\kangaekata 速さが一定の区間ごとに、①の長方形の面積（速さ×時間）を求めて、たし合わせます。\par
\begin{bkfigure}
\begin{tikzpicture}
\begin{axis}[bkplot, width=66mm, height=36mm, xmin=0, xmax=58, ymin=0, ymax=25, xtick={30,50}, ytick={10,20}, xlabel={$t$}, ylabel={$v$}]
  \fill[bkBlue!25] (axis cs:0,0) rectangle (axis cs:30,20);
  \fill[bkOrange!35] (axis cs:30,0) rectangle (axis cs:50,10);
  \addplot[bkcurve, line width=1.3pt] coordinates {(0,20) (30,20)};
  \addplot[bkcurve, line width=1.3pt] coordinates {(30,10) (50,10)};
  \node[font=\sffamily\footnotesize] at (axis cs:15,10) {600};
  \node[font=\sffamily\footnotesize] at (axis cs:40,5) {200};
\end{axis}
\end{tikzpicture}
\end{bkfigure}
\shiki はじめの30秒：$20\times30=600$\par
\hspace*{2.6\zw}次の20秒：$10\times20=200$\par
\hspace*{2.6\zw}合計：$600+200=800$\par
\kotae \qty{800}{m}
\end{reidai}

\begin{jibun}
自転車が、\qty{5}{m/s}で40秒走ったあと、\qty{3}{m/s}で60秒走りました。進んだ距離を求めましょう。
\bkjans{$5\times40+3\times60=200+180=\qty{380}{m}$}
\end{jibun}

\begin{reidai}[単位をそろえてから]
車が、時速\qty{90}{km}で2分走りました。進んだ距離は何mですか。\par
\kangaekata 速さは「1時間あたり」、時間は「分」で、単位がそろっていません。速さを m/s に、時間を s にそろえてから、かけます（2-10）。1 kmは1000 mなので \qty{1000}{m/km}、1時間は $60\times60=3600$ 秒なので \qty{3600}{s/h} を使います。\par
\shiki 速さ：$\qty{90}{km/h}\times\qty{1000}{m/km}=\qty{90000}{m/h}$\par
\hspace*{2.6\zw}$\qty{90000}{m/h}\div\qty{3600}{s/h}=\qty{25}{m/s}$\par
\hspace*{2.6\zw}時間：$\qty{2}{min}\times\qty{60}{s/min}=\qty{120}{s}$\par
\hspace*{2.6\zw}距離：$\qty{25}{m/s}\times\qty{120}{s}=\qty{3000}{m}$\par
\hspace*{2.6\zw}単位の約分：$90\,\dfrac{\cancel{\mathrm{km}}}{\mathrm{h}}\times1000\,\dfrac{\mathrm{m}}{\cancel{\mathrm{km}}}\div3600\,\dfrac{\mathrm{s}}{\mathrm{h}}=25\,\dfrac{\mathrm{m}}{\mathrm{s}}$\par
\hspace*{2.6\zw}確かめ（時間を h にそろえる）：$\qty{2}{min}\div\qty{60}{min/h}=\dfrac{1}{30}$\,h\par
\hspace*{2.6\zw}$\qty{90}{km/h}\times\dfrac{1}{30}\,\mathrm{h}=\qty{3}{km}$（＝\qty{3000}{m}）\par
\kotae \qty{3000}{m}（\qty{3}{km}）
\end{reidai}

\begin{jibun}
自転車が、時速\qty{18}{km}で3分走りました。進んだ距離は何mですか。
\bkjans{$\qty{18}{km/h}\times\qty{1000}{m/km}\div\qty{3600}{s/h}=\qty{5}{m/s}$、$\qty{3}{min}\times\qty{60}{s/min}=\qty{180}{s}$。$5\times180=\qty{900}{m}$}
\end{jibun}

\begin{reidai}[階段の和ではさむ]
止まっていた車が走り出し、速さが $v=3t$（m/s）でふえていきます。2秒間を0.5秒ずつの区間に分け、（1）はじめの速さを使った和、（2）おわりの速さを使った和を求め、（3）三角形の面積（ほんとうの距離）とくらべましょう。\par
\kangaekata ②・③と同じく、分けて（0.5秒ずつ4区間）、かけて（速さ×0.5）、たします。区間のはしの時刻は0、0.5、1、1.5、2秒で、そのときの速さは $3t$ に入れて0、1.5、3、4.5、6です。\par
\shiki（1）はじめの速さ：$(0+1.5+3+4.5)\times0.5=9\times0.5=4.5$\par
\hspace*{2.6\zw}（2）おわりの速さ：$(1.5+3+4.5+6)\times0.5=15\times0.5=7.5$\par
\hspace*{2.6\zw}（3）2秒後の速さは6。三角形：$2\times6\div2=6$\par
\begin{bkfigure}
\begin{tikzpicture}
\begin{axis}[bkplot, name=s1, width=48mm, height=38mm, xmin=0, xmax=2.4, ymin=0, ymax=7.2, xtick={0.5,1,1.5,2}, ytick={3,6}, xlabel={$t$}, ylabel={$v$}, x tick label style={font=\sffamily\scriptsize}, title={\sffamily\footnotesize（1）はじめの速さ}]
  \foreach \k in {1,2,3} {\edef\tmp{\noexpand\fill[bkBlue!25, draw=bkBlue!70, line width=0.4pt] (axis cs:{0.5*\k},0) rectangle (axis cs:{0.5*\k+0.5},{1.5*\k});}\tmp}
  \addplot[bkcurve, line width=1.2pt, domain=0:2.1] {3*x};
\end{axis}
\begin{axis}[bkplot, at={(s1.east)}, anchor=west, xshift=12mm, width=48mm, height=38mm, xmin=0, xmax=2.4, ymin=0, ymax=7.2, xtick={0.5,1,1.5,2}, ytick={3,6}, xlabel={$t$}, ylabel={$v$}, x tick label style={font=\sffamily\scriptsize}, title={\sffamily\footnotesize（2）おわりの速さ}]
  \foreach \k in {0,1,2,3} {\edef\tmp{\noexpand\fill[bkBlue!25, draw=bkBlue!70, line width=0.4pt] (axis cs:{0.5*\k},0) rectangle (axis cs:{0.5*\k+0.5},{1.5*\k+1.5});}\tmp}
  \addplot[bkcurve, line width=1.2pt, domain=0:2.1] {3*x};
\end{axis}
\end{tikzpicture}
\end{bkfigure}
\kotae（1）\qty{4.5}{m}\quad（2）\qty{7.5}{m}\quad（3）\qty{6}{m}。$4.5<6<7.5$ で、ほんとうの距離は2つの和のあいだにある
\end{reidai}

\begin{jibun}
止まっていた車の速さが $v=4t$（m/s）でふえていきます。4秒間を1秒ずつの区間に分け、（1）はじめの速さを使った和、（2）おわりの速さを使った和を求め、（3）三角形の面積とくらべましょう。
\bkjans{（1）\nobreak $(0+4+8+12)\times1=\qty{24}{m}$\quad（2）\nobreak $(4+8+12+16)\times1=\qty{40}{m}$\quad（3）\nobreak $4\times16\div2=\qty{32}{m}$。$24<32<40$}
\end{jibun}

\begin{reidai}[止まっていた物が速くなる（三角形）]
止まっていた車が、4秒で\qty{12}{m/s}まで一定の割合で速くなりました。この4秒で進んだ距離を求めましょう。\par
\kangaekata $v$-$t$ グラフは、原点（0秒で速さ0）から（4，12）までの直線。下は、底辺4（時間）、高さ12（4秒後の速さ）の三角形です（④）。\par
\shiki 底辺4、高さ12の三角形：$4\times12\div2=24$\par
\hspace*{2.6\zw}確かめ（平均の速さ）：$(0+12)\div2=6$、$6\times4=24$\par
\begin{bkfigure}
\begin{tikzpicture}
\begin{axis}[bkplot, width=50mm, height=36mm, xmin=0, xmax=5, ymin=0, ymax=14.5, xtick={4}, ytick={12}, xlabel={$t$}, ylabel={$v$}]
  \fill[bkBlue!25] (axis cs:0,0) -- (axis cs:4,12) -- (axis cs:4,0) -- cycle;
  \draw[bkGray, dashed, line width=0.5pt] (axis cs:0,12) -- (axis cs:4,12);
  \addplot[bkcurve, line width=1.3pt, domain=0:4] {3*x};
  \node[font=\sffamily\footnotesize] at (axis cs:2.9,3) {24};
\end{axis}
\end{tikzpicture}
\end{bkfigure}
\kotae \qty{24}{m}
\end{reidai}

\begin{jibun}
止まっていた電車が、20秒で\qty{15}{m/s}まで一定の割合で速くなりました。この20秒で進んだ距離を求めましょう。
\bkjans{$20\times15\div2=\qty{150}{m}$}
\end{jibun}

\begin{reidai}[加速度がわかっているとき]
止まっていた物体が、加速度\qty{2}{m/s^2}で5秒間速くなりました。進んだ距離を求めましょう。\par
\noindent\begin{minipage}[c]{0.6\linewidth}\raggedright
\kangaekata まず、5秒後の速さ（三角形の高さ）を求めます。加速度\qty{2}{m/s^2}は「1秒ごとに\qty{2}{m/s}ずつ速くなる」なので、5秒で $2\times5$ ふえます。\par
\shiki 5秒後の速さ：$2\times5=\qty{10}{m/s}$\par
\hspace*{2.6\zw}三角形の面積：$5\times10\div2=25$\par
\hspace*{2.6\zw}確かめ（$\frac{1}{2}at^2$）：$\dfrac{1}{2}\times2\times5^2=25$\par
\end{minipage}\hfill
\begin{minipage}[c]{0.37\linewidth}\centering
\begin{tikzpicture}
\begin{axis}[bkplot, width=41mm, height=38mm, xmin=0, xmax=6.2, ymin=0, ymax=12.5, xtick={1,...,5}, ytick={2,4,6,8,10}, xlabel={$t$}, ylabel={$v$}, y tick label style={font=\sffamily\scriptsize}]
  \fill[bkBlue!25] (axis cs:0,0) -- (axis cs:5,10) -- (axis cs:5,0) -- cycle;
  \addplot[bkcurve, line width=1.3pt, domain=0:5] {2*x};
  \node[font=\sffamily\footnotesize] at (axis cs:3.7,2.4) {25};
\end{axis}
\end{tikzpicture}
\end{minipage}\par
\kotae \qty{25}{m}
\end{reidai}

\begin{jibun}
止まっていた物体が、加速度\qty{1.5}{m/s^2}で4秒間速くなりました。進んだ距離を求めましょう。
\bkjans{4秒後の速さ $1.5\times4=\qty{6}{m/s}$。$4\times6\div2=\qty{12}{m}$}
\end{jibun}

\begin{reidai}[はじめから動いている（台形）]
\qty{5}{m/s}で走っていた自転車が、一定の割合で速くなり、4秒後に\qty{15}{m/s}になりました。この4秒で進んだ距離を求めましょう。\par
\noindent\begin{minipage}[c]{0.6\linewidth}\raggedright
\kangaekata ⑤の台形です。たての辺（上底・下底）が5と15、横のはば（高さ）が4です。\par
\shiki $(5+15)\times4\div2=20\times4\div2=40$\par
\hspace*{2.6\zw}確かめ（長方形＋三角形）：$5\times4+4\times(15-5)\div2=20+20=40$\par
\end{minipage}\hfill
\begin{minipage}[c]{0.37\linewidth}\centering
\begin{tikzpicture}
\begin{axis}[bkplot, width=41mm, height=38mm, xmin=0, xmax=5, ymin=0, ymax=17.5, xtick={4}, ytick={5,15}, xlabel={$t$}, ylabel={$v$}]
  \fill[bkBlue!25] (axis cs:0,0) -- (axis cs:0,5) -- (axis cs:4,15) -- (axis cs:4,0) -- cycle;
  \draw[bkGray, dashed, line width=0.5pt] (axis cs:4,0) -- (axis cs:4,15) (axis cs:0,15) -- (axis cs:4,15);
  \addplot[bkcurve, line width=1.3pt, domain=0:4] {5+2.5*x};
  \node[font=\sffamily\footnotesize] at (axis cs:2,3.5) {40};
\end{axis}
\end{tikzpicture}
\end{minipage}\par
\kotae \qty{40}{m}
\end{reidai}

\begin{jibun}
\qty{4}{m/s}で走っていた自転車が、一定の割合で速くなり、3秒後に\qty{10}{m/s}になりました。この3秒で進んだ距離を求めましょう。
\bkjans{$(4+10)\times3\div2=\qty{21}{m}$}
\end{jibun}

\begin{reidai}[ブレーキでおそくなる]
（1）\qty{15}{m/s}で走っていた車が、一定の割合でおそくなり、4秒後に\qty{5}{m/s}になりました。この4秒で進んだ距離は？\par
（2）時速\qty{72}{km}で走っていた車が、加速度 $\qty{-4}{m/s^2}$ でブレーキをかけました。止まるまでに進む距離は？\par
\kangaekata（1）は右下がりの台形（⑥）。（2）は、速さを m/s にしてから、止まるまでの時間（三角形の底辺）を求めます。加速度 $-4$ は「1秒ごとに\qty{4}{m/s}ずつおそくなる」です。\par
\begin{bkfigure}
\begin{tikzpicture}
\begin{axis}[bkplot, name=e1, width=46mm, height=36mm, xmin=0, xmax=5, ymin=0, ymax=17.5, xtick={4}, ytick={5,15}, xlabel={$t$}, ylabel={$v$}, title={\sffamily\footnotesize（1）}]
  \fill[bkBlue!25] (axis cs:0,0) -- (axis cs:0,15) -- (axis cs:4,5) -- (axis cs:4,0) -- cycle;
  \draw[bkGray, dashed, line width=0.5pt] (axis cs:4,0) -- (axis cs:4,5);
  \addplot[bkcurve, line width=1.3pt, domain=0:4] {15-2.5*x};
  \node[font=\sffamily\footnotesize] at (axis cs:2,3.6) {40};
\end{axis}
\begin{axis}[bkplot, at={(e1.east)}, anchor=west, xshift=12mm, width=46mm, height=36mm, xmin=0, xmax=6.2, ymin=0, ymax=23, xtick={5}, ytick={20}, xlabel={$t$}, ylabel={$v$}, title={\sffamily\footnotesize（2）}]
  \fill[bkBlue!25] (axis cs:0,0) -- (axis cs:0,20) -- (axis cs:5,0) -- cycle;
  \addplot[bkcurve, line width=1.3pt, domain=0:5] {20-4*x};
  \node[font=\sffamily\footnotesize] at (axis cs:1.3,5) {50};
\end{axis}
\end{tikzpicture}
\end{bkfigure}
\shiki（1）$(15+5)\times4\div2=40$\par
\hspace*{2.6\zw}（2）$\qty{72}{km/h}\times\qty{1000}{m/km}\div\qty{3600}{s/h}=\qty{20}{m/s}$\par
\hspace*{4.6\zw}止まるまでの時間：$20\div4=\qty{5}{s}$\par
\hspace*{4.6\zw}三角形の面積：$5\times20\div2=50$\par
\kotae（1）\qty{40}{m}\quad（2）\qty{50}{m}
\end{reidai}

\begin{jibun}
（1）\qty{12}{m/s}で走っていた車が、一定の割合でおそくなり、2秒後に\qty{6}{m/s}になりました。この2秒で進んだ距離は？\par
（2）時速\qty{54}{km}で走っていた車が、加速度 $\qty{-3}{m/s^2}$ でブレーキをかけました。止まるまでに進む距離は？
\bkjans{（1）\nobreak $(12+6)\times2\div2=\qty{18}{m}$\quad（2）\nobreak \qty{15}{m/s}、$15\div3=\qty{5}{s}$、$5\times15\div2=\qty{37.5}{m}$}
\end{jibun}

\begin{reidai}[3つの区間：駅から駅まで]
電車が駅を出て、20秒で\qty{15}{m/s}まで一定の割合で速くなり、\qty{15}{m/s}で60秒走り、10秒で一定の割合でおそくなって次の駅に止まりました。駅と駅のあいだの距離を求めましょう。\par
\kangaekata $v$-$t$ グラフは、三角形・長方形・三角形をつなげた台形です。区間ごとに求めてたします。\par
\begin{bkfigure}
\begin{tikzpicture}
\begin{axis}[bkplot, width=84mm, height=36mm, xmin=0, xmax=98, ymin=0, ymax=18.5, xtick={20,80,90}, ytick={15}, xlabel={$t$}, ylabel={$v$}]
  \fill[bkYellow!45] (axis cs:0,0) -- (axis cs:20,15) -- (axis cs:20,0) -- cycle;
  \fill[bkBlue!22] (axis cs:20,0) rectangle (axis cs:80,15);
  \fill[bkYellow!45] (axis cs:80,0) -- (axis cs:80,15) -- (axis cs:90,0) -- cycle;
  \addplot[bkcurve, line width=1.3pt] coordinates {(0,0) (20,15) (80,15) (90,0)};
  \node[font=\sffamily\footnotesize] at (axis cs:14.5,4) {150};
  \node[font=\sffamily\footnotesize] at (axis cs:50,7.5) {900};
  \node[font=\sffamily\footnotesize] at (axis cs:83.4,4) {75};
\end{axis}
\end{tikzpicture}
\end{bkfigure}
\shiki 速くなる20秒：$20\times15\div2=150$\par
\hspace*{2.6\zw}一定の60秒：$15\times60=900$\par
\hspace*{2.6\zw}おそくなる10秒：$10\times15\div2=75$\par
\hspace*{2.6\zw}合計：$150+900+75=1125$\par
\hspace*{2.6\zw}確かめ（大きな台形：上底60、下底90、高さ15）：\par
\hspace*{2.6\zw}$(60+90)\times15\div2=1125$\par
\kotae \qty{1125}{m}
\end{reidai}

\begin{jibun}
電車が、10秒で\qty{10}{m/s}まで一定の割合で速くなり、\qty{10}{m/s}で30秒走り、10秒で一定の割合でおそくなって止まりました。進んだ距離を求めましょう。
\bkjans{$50+300+50=\qty{400}{m}$}
\end{jibun}

\begin{reidai}[マイナスの面積]
エレベーターが、\qty{2}{m/s}で6秒上がり、そのあと\qty{-3}{m/s}で2秒下がりました（上向きが＋）。変位と、動いた道のりを求めましょう。\par
\noindent\begin{minipage}[c]{0.6\linewidth}\raggedright
\kangaekata 下がるときの速度はマイナスなので、⑦のマイナスの面積です。変位は符号のままたし、道のりは大きさ（プラス）をたします。\par
\shiki 上がる：$2\times6=12$、下がる：$(-3)\times2=-6$\par
\hspace*{2.6\zw}変位：$12+(-6)=6$\quad 道のり：$12+6=18$\par
\end{minipage}\hfill
\begin{minipage}[c]{0.37\linewidth}\centering
\begin{tikzpicture}
\begin{axis}[bkplot, width=41mm, height=40mm, xmin=0, xmax=9.6, ymin=-4.6, ymax=3.4, xtick={6}, ytick={-3,2}, xlabel={$t$}, ylabel={$v$}]
  \fill[bkBlue!22] (axis cs:0,0) rectangle (axis cs:6,2);
  \fill[bkRed!25] (axis cs:6,0) rectangle (axis cs:8,-3);
  \addplot[bkcurve, line width=1.3pt] coordinates {(0,2) (6,2)};
  \addplot[bkcurve, line width=1.3pt] coordinates {(6,-3) (8,-3)};
  \draw[bkGray, line width=0.4pt] (axis cs:8,0) -- (axis cs:8,0.5);
  \node[font=\sffamily\footnotesize, anchor=south] at (axis cs:8,0.5) {8};
  \node[font=\sffamily\footnotesize] at (axis cs:3,1) {$+12$};
  \node[font=\sffamily\footnotesize] at (axis cs:7,-1.5) {$-6$};
\end{axis}
\end{tikzpicture}
\end{minipage}\par
\kotae 変位\qty{6}{m}（はじめより上）、道のり\qty{18}{m}
\end{reidai}

\begin{jibun}
ボールが\qty{3}{m/s}で2秒ころがり、かべではね返って\qty{-2}{m/s}で4秒ころがりました。変位と道のりを求めましょう。
\bkjans{変位 $6+(-8)=\qty{-2}{m}$（はじめの位置より\qty{2}{m}手前）、道のり $6+8=\qty{14}{m}$}
\end{jibun}

\begin{reidai}[線が横軸を横切る：投げ上げ]
ビルの屋上から、ボールを\qty{10}{m/s}で真上に投げ上げました（上向きが＋、加速度 $\qty{-10}{m/s^2}$）。速度は $v=10-10t$ です。\par
（1）いちばん高くなるまでに、何m上がりますか。（2）3秒後の変位はいくつですか。\par
\noindent\begin{minipage}[c]{0.6\linewidth}\raggedright
\kangaekata $v=0$ になる時刻で、上の三角形と下の三角形に分けます。$10-10t=0$ より $t=1$。3秒後の速度は $10-30=-20$。\par
\shiki（1）0〜1秒の三角形：$1\times10\div2=5$\par
\hspace*{2.6\zw}（2）1〜3秒の三角形：$-(2\times20\div2)=-20$\par
\hspace*{4.6\zw}変位：$5+(-20)=-15$\par
\end{minipage}\hfill
\begin{minipage}[c]{0.37\linewidth}\centering
\begin{tikzpicture}
\begin{axis}[bkplot, width=41mm, height=44mm, xmin=0, xmax=3.7, ymin=-23, ymax=13, xtick={2}, ytick={-20,10}, xlabel={$t$}, ylabel={$v$}]
  \fill[bkBlue!25] (axis cs:0,0) -- (axis cs:0,10) -- (axis cs:1,0) -- cycle;
  \fill[bkRed!25] (axis cs:1,0) -- (axis cs:3,-20) -- (axis cs:3,0) -- cycle;
  \draw[bkGray, line width=0.4pt] (axis cs:1,0) -- (axis cs:1,1) (axis cs:3,0) -- (axis cs:3,1);
  \addplot[bkcurve, line width=1.3pt, domain=0:3.2] {10-10*x};
  \node[font=\sffamily\footnotesize, anchor=south] at (axis cs:1,1) {1};
  \node[font=\sffamily\footnotesize, anchor=south] at (axis cs:3,1) {3};
  \node[font=\sffamily\scriptsize] at (axis cs:0.3,2.6) {$+5$};
  \node[font=\sffamily\footnotesize] at (axis cs:2.55,-6) {$-20$};
\end{axis}
\end{tikzpicture}
\end{minipage}\par
\kotae（1）\qty{5}{m}\quad（2）\qty{-15}{m}（投げた所より\qty{15}{m}下）
\end{reidai}

\begin{jibun}
ボールを\qty{15}{m/s}で真上に投げ上げました（上向きが＋、加速度 $\qty{-10}{m/s^2}$）。速度は $v=15-10t$ です。（1）いちばん高くなるまでに何m上がりますか。（2）3秒後の変位はいくつですか。
\bkjans{（1）\nobreak $t=1.5$、$1.5\times15\div2=\qty{11.25}{m}$\quad（2）\nobreak 3秒後の速度 $-15$。$11.25+(-11.25)=\qty{0}{m}$（投げた所にもどる）}
\end{jibun}

\begin{reidai}[面積をつみ上げる]
物体の速さを、1秒ごとに記録しました（それぞれのあいだは一定の割合で変わる）。0秒の位置を0として、1秒ごとの位置を求めましょう。\par
\begin{bkcenter}
\begin{bktable}{colspec={Q[l,m,wd=22mm] *{4}{Q[c,m,wd=10mm]}}, rowsep=1pt}
$t$（s） & 0 & 1 & 2 & 3 \\
$v$（m/s） & 2 & 4 & 4 & 0 \\
\end{bktable}
\end{bkcenter}
\noindent\begin{minipage}[c]{0.6\linewidth}\raggedright
\kangaekata ⑧のように、1秒ごとの面積（台形・長方形・三角形）を求め、前の位置につみ上げます。\par
\shiki 0〜1秒（台形）：$(2+4)\times1\div2=3$\par
\hspace*{2.6\zw}1〜2秒（長方形）：$4\times1=4$\par
\hspace*{2.6\zw}2〜3秒（三角形）：$1\times4\div2=2$\par
\hspace*{2.6\zw}位置：$0+3=3$、$3+4=7$、$7+2=9$\par
\end{minipage}\hfill
\begin{minipage}[c]{0.37\linewidth}\centering
\begin{tikzpicture}
\begin{axis}[bkplot, width=41mm, height=38mm, xmin=0, xmax=3.7, ymin=0, ymax=5.3, xtick={1,2,3}, ytick={2,4}, xlabel={$t$}, ylabel={$v$}]
  \fill[bkYellow!45] (axis cs:0,0) -- (axis cs:0,2) -- (axis cs:1,4) -- (axis cs:1,0) -- cycle;
  \fill[bkBlue!22] (axis cs:1,0) rectangle (axis cs:2,4);
  \fill[bkGreen!25] (axis cs:2,0) -- (axis cs:2,4) -- (axis cs:3,0) -- cycle;
  \addplot[bkcurve, line width=1.3pt] coordinates {(0,2) (1,4) (2,4) (3,0)};
  \node[font=\sffamily\footnotesize] at (axis cs:0.5,1.2) {3};
  \node[font=\sffamily\footnotesize] at (axis cs:1.5,2) {4};
  \node[font=\sffamily\footnotesize] at (axis cs:2.3,1) {2};
\end{axis}
\end{tikzpicture}
\end{minipage}\par
\kotae 1秒で\qty{3}{m}、2秒で\qty{7}{m}、3秒で\qty{9}{m}
\end{reidai}

\begin{jibun}
速さが、0秒で0、2秒で\qty{6}{m/s}、4秒で\qty{6}{m/s}、6秒で\qty{2}{m/s}（あいだは一定の割合で変わる）の物体の、2秒ごとの位置を求めましょう。
\bkjans{\mbox{2秒}で\qty{6}{m}、\mbox{4秒}で\qty{18}{m}、\mbox{6秒}で\qty{26}{m}}
\end{jibun}

%% ---------------- まちがえやすいところ・物理工学 ----------------
\bksection{tsumazuki}{まちがえやすいところをチェック}

\begin{tsumazuki}{単位をそろえずに面積を計算する}
\bkng{時速\qty{72}{km}で10秒：$72\times10=720$ で\qty{720}{m}}
\bkok{$\qty{72}{km/h}\times\qty{1000}{m/km}\div\qty{3600}{s/h}=\qty{20}{m/s}$。$20\times10=\qty{200}{m}$}
「1時間あたりの速さ」に「秒」をかけても、意味のある数になりません。かける前に、単位を m/s と s（または km/h と h）にそろえます。面積の単位が、距離の単位になっているかを確かめましょう。
\end{tsumazuki}

\begin{tsumazuki}{三角形の「÷2」をわすれる}
\bkng{止まっていた車が4秒で\qty{12}{m/s}：$12\times4=\qty{48}{m}$}
\bkok{$4\times12\div2=\qty{24}{m}$}
\noindent\begin{minipage}[c]{0.62\linewidth}\raggedright
$12\times4$ は「はじめから\qty{12}{m/s}で4秒走った」ときの距離で、右の図の長方形全体です。ほんとうは0からだんだん速くなるので、その半分の三角形（青）です。平均の速さ（$12\div2=6$）で確かめるとまちがえません。
\end{minipage}\hfill
\begin{minipage}[c]{0.34\linewidth}\centering
\begin{tikzpicture}
\begin{axis}[bkplot, width=40mm, height=32mm, xmin=0, xmax=5, ymin=0, ymax=14.5, xtick={4}, ytick={12}, xlabel={$t$}, ylabel={$v$}]
  \fill[bkGray!20] (axis cs:0,0) rectangle (axis cs:4,12);
  \fill[bkBlue!30] (axis cs:0,0) -- (axis cs:4,12) -- (axis cs:4,0) -- cycle;
  \addplot[bkcurve, line width=1.2pt, domain=0:4] {3*x};
\end{axis}
\end{tikzpicture}
\end{minipage}
\end{tsumazuki}

\begin{tsumazuki}{変位と道のりをまちがえる}
\bkng{\qty{4}{m/s}で3秒、\qty{-2}{m/s}で2秒。道のりは $12+(-4)=\qty{8}{m}$}
\bkok{道のりは $12+4=\qty{16}{m}$。変位が\qty{8}{m}}
もどった分も、実際に動いた長さです。道のりはマイナスの面積もプラスにしてたし、変位は符号のままたします。
\end{tsumazuki}

\begin{tsumazuki}{$x$-$t$ グラフの面積を求めてしまう}
距離になるのは、$v$-$t$ グラフ（速度のグラフ）の面積です。$x$-$t$ グラフ（位置のグラフ）では、ある時刻の位置は、グラフの高さそのものです。たて軸が何かを、いつも確かめましょう。
\end{tsumazuki}

\begin{column}{コンピュータは階段で計算する}
車の走行距離計や、スマートフォンの歩きの記録アプリは、0.1秒ごとなどの短い時間ごとに速さをはかり、「速さ×時間」をつぎつぎにたして、進んだ距離を出しています。この章の「細かく分けて、かけて、たす」そのものです。区間が短いので、階段とほんとうのグラフのすき間はとても小さく、正確な距離が求められます。5-8のオイラー法も、同じ計算でした。
\end{column}

\begin{butsuri}
\begin{itemize}[leftmargin=1.2\zw, itemsep=2pt, topsep=0pt, label={\color{bkPurple}●}]
\item \textbf{タコグラフ}\\
記録された速さのグラフの面積から、走った距離を計算します（はてな？）。運転時間や休けい時間も、グラフからわかります。
\item \textbf{電車の運転}\\
電車は、加速・一定の速さ・ブレーキを組み合わせて走ります。$v$-$t$ グラフの面積が駅と駅のあいだの距離になるように、運転のしかたを決めています（下の図）。
\item \textbf{ブレーキの距離}\\
同じ加速度でブレーキをかけると、速さが2倍なら止まるまでの時間も2倍です。三角形の底辺も高さも2倍になるので、面積（止まるまでの距離）は $2\times2=4$ 倍になります（下の右の図）。
\end{itemize}
\begin{bkfigure}
\begin{tikzpicture}
\begin{axis}[bkplot, name=r1, width=70mm, height=40mm, xmin=0, xmax=128, ymin=0, ymax=25, xtick={30,90,110}, ytick={20}, xlabel={$t$}, ylabel={$v$}]
  \fill[bkPurple!15] (axis cs:0,0) -- (axis cs:30,20) -- (axis cs:90,20) -- (axis cs:110,0) -- cycle;
  \addplot[bkcurve, line width=1.5pt] coordinates {(0,0) (30,20) (90,20) (110,0)};
  \node[font=\sffamily\scriptsize] at (axis cs:12,22) {加速};
  \node[font=\sffamily\scriptsize] at (axis cs:60,22.5) {一定の速さ};
  \node[font=\sffamily\scriptsize] at (axis cs:106,22) {ブレーキ};
  \node[font=\sffamily\scriptsize] at (axis cs:60,9) {面積＝駅と駅の距離};
\end{axis}
\begin{axis}[bkplot, at={(r1.east)}, anchor=west, xshift=12mm, width=40mm, height=40mm, xmin=0, xmax=4.8, ymin=0, ymax=24, xtick={2,4}, ytick={10,20}, xlabel={$t$}, ylabel={$v$}]
  \fill[bkPurple!15] (axis cs:0,0) -- (axis cs:0,20) -- (axis cs:4,0) -- cycle;
  \fill[bkOrange!45] (axis cs:0,0) -- (axis cs:0,10) -- (axis cs:2,0) -- cycle;
  \addplot[bkcurve, line width=1.3pt, domain=0:4] {20-5*x};
  \addplot[bkcurve2, line width=1.3pt, domain=0:2] {10-5*x};
\end{axis}
\end{tikzpicture}
\bkfigcaption{左：電車の運転曲線（ランカーブ）。右：速さ2倍（むらさき）は、距離4倍。}
\end{bkfigure}
\end{butsuri}

%% ---------------- 章末問題 ----------------
\bkendproblems
\bkproblemgroup{基本}{2}

\bkproblem 時速\qty{54}{km}で2分走りました。進んだ距離は何mですか。
\bkpans{$\qty{54}{km/h}\times\qty{1000}{m/km}\div\qty{3600}{s/h}=\qty{15}{m/s}$、$\qty{2}{min}\times\qty{60}{s/min}=\qty{120}{s}$。$15\times120=\qty{1800}{m}$}

\bkproblem 止まっていた物体が、一定の加速度\qty{3}{m/s^2}で6秒動きました。$v$-$t$ グラフの面積から、進んだ距離を求めましょう。
\bkpans{\mbox{6秒後}の速さ $3\times6=\qty{18}{m/s}$。$6\times18\div2=\qty{54}{m}$}

\bkproblem \qty{10}{m/s}の自転車が、一定の割合でおそくなり、5秒後に\qty{4}{m/s}になりました。この5秒で進んだ距離を求めましょう。
\bkpans{$(10+4)\times5\div2=\qty{35}{m}$}

\bkproblem 電車が、30秒で\qty{20}{m/s}まで一定の割合で速くなり、\qty{20}{m/s}で60秒走り、20秒で一定の割合でおそくなって止まりました。駅と駅のあいだの距離を求めましょう。
\bkpans{$300+1200+200=\qty{1700}{m}$}

\bkproblem \qty{4}{m/s}で5秒進み、そのあと\qty{-3}{m/s}で4秒もどりました。変位と道のりを求めましょう。
\bkpans{変位 $20+(-12)=\qty{8}{m}$、道のり $20+12=\qty{32}{m}$}

\bkproblem 速さが、0秒で0、1秒で\qty{2}{m/s}、2秒で\qty{2}{m/s}、3秒で\qty{4}{m/s}、4秒で0（あいだは一定の割合で変わる）の物体の、1秒ごとの位置を求めましょう（0秒の位置は0）。
\bkpans{1秒で\qty{1}{m}、2秒で\qty{3}{m}、3秒で\qty{6}{m}、4秒で\qty{8}{m}}

\bkproblemgroup{チャレンジ}{2}

\bkproblem ボールを\qty{20}{m/s}で真上に投げ上げます（上向きが＋、加速度 \qty{-10}{m/s^2}）。\par
（1）$v$-$t$ グラフをかきましょう。\quad（2）いちばん高くなるまでの距離を、面積で求めましょう。\quad（3）4秒後の変位はいくつですか。
\bkpans{（1）\nobreak $v=20-10t$ の直線（2秒で0）\quad（2）\nobreak $2\times20\div2=\qty{20}{m}$\quad（3）\nobreak $20+(-20)=\qty{0}{m}$（もとの高さにもどる）}

\bkproblem 時速\qty{36}{km}と時速\qty{72}{km}の車が、どちらも加速度 \qty{-5}{m/s^2} でブレーキをかけて止まりました。止まるまでの距離をそれぞれ求め、くらべましょう。
\bkpans{\qty{10}{m/s}：\mbox{2秒}で $2\times10\div2=\qty{10}{m}$。\qty{20}{m/s}：\mbox{4秒}で $4\times20\div2=\qty{40}{m}$。速さ\mbox{2倍}で距離は\mbox{4倍}}

\bkproblem 止まっていた物体の速さが $v=6t$（m/s）でふえていきます。2秒間を0.5秒ずつに分け、はじめの速さを使った和と、おわりの速さを使った和を求めましょう。ほんとうの距離はいくつで、2つの和とどんな関係にありますか。
\bkpans{はじめ $(0+3+6+9)\times0.5=\qty{9}{m}$、おわり $(3+6+9+12)\times0.5=\qty{15}{m}$。ほんとうは $2\times12\div2=\qty{12}{m}$ で、2つの和のあいだ}

\bkproblem バスが、0.25時間で時速\qty{40}{km}まで一定の割合で速くなり、時速\qty{40}{km}で1.5時間走り、0.25時間で一定の割合でおそくなって止まりました。進んだ距離は何kmですか。
\bkpans{加速と減速は三角形で、どちらも $0.25\times40\div2=5$。一定の速さの所は $40\times1.5=60$。合計 $5+60+5=\qty{70}{km}$}

\begin{furikaeri}
\begin{checklist}
\item 速さが変わるときも、細かく分けて、かけて、たして距離を求められる。\hfill{\footnotesize\sffamily\color{bkTheme}→\pageref{bk:5-9:wakeru}ページ}
\item 長方形・三角形・台形の面積で、進んだ距離がわかる。\hfill{\footnotesize\sffamily\color{bkTheme}\mbox{→\pageref{bk:5-9:sankaku}ページ}}
\item マイナスの面積の意味（変位と道のり）を説明できる。\hfill{\footnotesize\sffamily\color{bkTheme}→\pageref{bk:5-9:minus}ページ}
\end{checklist}
\end{furikaeri}

\bkfukidashi{mina}{グラフがまっすぐなら、三角形や台形で計算できたね。じゃあ、グラフが曲がっていたら？　階段の考え方が、そのまま使えるのかな？}
